\documentclass[10pt,aspectratio=169]{beamer}
\usetheme{Madrid}

% --- Pacotes Recomendados para Overleaf ---
\usepackage[utf8]{inputenc}
\usepackage[portuguese]{babel}
\usepackage{amsmath}
\usepackage{amsfonts}
\usepackage{amssymb}
\usepackage{graphicx}
\usepackage{booktabs}
\usepackage{tikz}
\usepackage{hyperref}

% --- Customização de Cores (Tema Escuro/Tecnológico Elegante) ---
\definecolor{DeepSlate}{RGB}{15, 23, 42}
\definecolor{TechBlue}{RGB}{14, 165, 233}
\definecolor{LightGray}{RGB}{241, 245, 249}
\definecolor{DarkText}{RGB}{33, 37, 41}

\setbeamercolor{palette primary}{bg=DeepSlate,fg=white}
\setbeamercolor{palette secondary}{bg=TechBlue,fg=white}
\setbeamercolor{palette tertiary}{bg=DeepSlate,fg=white}
\setbeamercolor{structure}{fg=TechBlue}
\setbeamercolor{titlelike}{parent=palette primary}
\setbeamercolor{background canvas}{bg=white}
\setbeamercolor{normal text}{fg=DarkText}

% --- Configurações Gerais ---
\title[Latent Diffusion com EMNIST]{Latent Diffusion com o Dataset EMNIST}
\subtitle{Implementação de um Pipeline Generativo Bidimensional (VAE + Denoiser)}

\author{%
  Arthur Jorge - 2022055718 \\
  Jose Adonai Arruda - 2026685902 \\
  Lucas Ribeiro \\
  Wagner
}

\institute{Ambiente de Desenvolvimento: PyTorch \& Google Colab}
\date{\today}

% --- REMOVER RODAPÉ E ÍCONES DE NAVEGAÇÃO ---
\setbeamertemplate{footline}{} 
\setbeamertemplate{navigation symbols}{}

\begin{document}

% ==============================================================================
% SLIDE 1: Capa
% ==============================================================================
\begin{frame}
    \titlepage
\end{frame}

% ==============================================================================
% SLIDE 2: Agenda
% ==============================================================================
\begin{frame}{Agenda}
    \tableofcontents
\end{frame}

\section{1. Visão Geral do Projeto}
% ==============================================================================
% SLIDE 3: Visão Geral do Pipeline
% ==============================================================================
\begin{frame}{Visão Geral do Projeto}
    O projeto implementa um modelo de \textbf{Latent Diffusion} aplicado ao dataset \textbf{EMNIST}. O pipeline divide-se em duas etapas macro fundamentais:
    \vspace{0.3cm}
    \begin{block}{1. Autoencoder Variacional (VAE)}
        Codifica as imagens originais de alta dimensionalidade para um espaço latente compacto, contínuo e denso ($m=32$), forçando o aprendizado de características estruturais robustas.
    \end{block}
    \begin{block}{2. Modelo de Difusão (Latent Diffusion)}
        Atua diretamente sobre o espaço latente reduzido. Um modelo \textit{Denoiser} aprende a reverter o processo de difusão determinística. O Decoder do VAE mapeia os vetores gerados de volta para o formato visual.
    \end{block}
\end{frame}

% ==============================================================================
% SLIDE 4: O Dataset EMNIST
% ==============================================================================
\begin{frame}{O Desafio do Dataset EMNIST}
    \begin{itemize}
        \item \textbf{Complexidade Superior:} Ao contrário do MNIST tradicional (apenas 10 dígitos), o EMNIST inclui dígitos e letras maiúsculas e minúsculas (A-Z, a-z, 0-9).
        \item \textbf{Variabilidade Alta:} Maior diversidade de traços manuais, rotações e distorções morfológicas.
        \item \textbf{Abordagem Numérica:} Imagens em tons de cinza normalizadas, tratadas como matrizes binárias ou contínuas dependendo do modelo de verossimilhança adotado.
    \end{itemize}
    \vspace{0.3cm}
    \begin{exampleblock}{Ambiente de Execução}
        Desenvolvido em \textbf{PyTorch} com aceleração por GPU no ambiente \textbf{Google Colab}, permitindo rápida prototipagem e tuning de hiperparâmetros.
    \end{exampleblock}
\end{frame}

\section{2. O Autoencoder Variacional (VAE)}
% ==============================================================================
% SLIDE 5: Arquitetura do VAE
% ==============================================================================
\begin{frame}{Arquitetura do VAE e Seleção de Blocos}
    Durante a modelagem no \texttt{VAE\_EMNIST\_tunning.ipynb}, testou-se o impacto estrutural da rede convolucional:
    \vspace{0.2cm}
    \begin{columns}
        \begin{column}{0.5\textwidth}
            \textbf{Configuração A: 2 Blocos Convolucionais}
            \begin{itemize}
                \item \texttt{Encoder2} e \texttt{Decoder2}.
                \item Menor profundidade, extração de features mais superficiais.
                \item Reconstruções ligeiramente borradas.
            \end{itemize}
        \end{column}
        \begin{column}{0.5\textwidth}
            \textbf{Configuração B: 3 Blocos Convolucionais (Vencedor)}
            \begin{itemize}
                \item \texttt{Encoder3} e \texttt{Decoder3}.
                \item Compressão hierárquica mais eficiente.
                \item Mantém alta fidelidade estrutural no gargalo.
            \end{itemize}
        \end{column}
    \end{columns}
    \vspace{0.4cm}
    \textbf{Dimensão Latente Ideal ($m=32$):} Provou ser o balanço perfeito entre compressão (gargalo restrito para a difusão) e capacidade de reconstrução da imagem.
\end{frame}

% ==============================================================================
% SLIDE 6: Escolha do Modelo de Likelihood
% ==============================================================================
\begin{frame}{Escolha do Modelo de Likelihood (Verossimilhança)}
    A escolha de $\tilde{\rho}_{X|Z}(x, z)$ dita a formulação matemática do termo de reconstrução $R$ (log-verossimilhança negativa):
    \vspace{0.3cm}
    \begin{itemize}
        \item \textbf{Gaussiana (MSE Loss):} Simétrica e derivável em toda parte. Tende a gerar imagens excessivamente borradas devido à penalidade quadrática que suaviza os gradientes.
        \item \textbf{Bernoulli (BCE Loss) -- Escolhida:} Melhor alinhamento com a natureza binária do EMNIST. Como a perda cresce infinitamente ao errar os extremos ($0$ e $1$), força reconstruções muito mais nítidas nas bordas das letras e dígitos.
        \item \textbf{Laplace (L1 Loss):} Robusta a outliers e preserva melhor as bordas que o MSE, porém mostra-se inferior ao BCE para o domínio específico do EMNIST.
    \end{itemize}
\end{frame}

% ==============================================================================
% SLIDE 7: Curvas de Custo por Pixel
% ==============================================================================
\begin{frame}{Análise Estatística: Curvas de Custo por Pixel}
    \begin{itemize}
        \item Comparação do custo de reconstrução em função do valor predito pelo Decoder para um pixel alvo fixo.
        \item As três curvas divergem agressivamente perto dos extremos ($0$ e $1$), região onde se concentra a vasta maioria dos pixels do EMNIST.
    \end{itemize}
    
    \begin{center}
        % ------------------------------------------------------------------
        % IMAGEM: likelihood_cost_curves.png
        % ------------------------------------------------------------------
        \includegraphics[width=0.75\textwidth]{likelihood_cost_curves.png}
        \\ \vspace{0.2cm}
        \footnotesize \textbf{Figura 1:} Custo de reconstrução vs. predição do decoder para cada função de verossimilhança.
    \end{center}
    
    \vfill
\end{frame}

% ==============================================================================
% SLIDE 8: Posterior Collapse e KL Annealing
% ==============================================================================
\begin{frame}{O Problema do Posterior Collapse e o KL Annealing}
    A função de perda do VAE equilibra duas forças opostas: $\text{Loss} = R + \beta L$. Onde $R$ é a reconstrução e $L$ é a divergência KL.
    \vspace{0.2cm}
    \begin{itemize}
        \item \textbf{Posterior Collapse (Colapso Posterior):} No início do treino, o Decoder ainda é ineficiente. O otimizador escolhe o "atalho preguiçoso": zera o termo KL ($L \to 0$), fazendo o Encoder prever sempre a distribuição padrão $\mathcal{N}(0, I)$ ignorando a imagem de entrada. O espaço latente colapsa em ruído inútil.
        \item \textbf{A Solução (KL Annealing):} Introduzir o peso dinâmico $\beta$ que cresce gradualmente ao longo das épocas. Permite que a rede aprenda a reconstruir ($R$) primeiro, para só então regularizar o espaço latente ($L$).
    \end{itemize}
\end{frame}

% ==============================================================================
% SLIDE 9: Implementação do KL Annealing
% ==============================================================================
\begin{frame}{Cronograma do KL Annealing (Schedule em Degraus)}
    O projeto adota um schedule em degraus, onde $\beta$ se eleva em $+0.2$ a cada 5 épocas, estabilizando em $1.0$.
    \vspace{0.3cm}
    \begin{table}[]
        \centering
        \begin{tabular}{lcccccc}
            \toprule
            \textbf{Época} & 1--4 & 5--9 & 10--14 & 15--19 & 20--24 & $\ge 25$ \\
            \midrule
            $\boldsymbol{\beta}$ & 0.0 & 0.2 & 0.4 & 0.6 & 0.8 & 1.0 \\
            \bottomrule
        \end{tabular}
    \end{table}
    \vspace{0.2cm}
    \begin{itemize}
        \item \textbf{Avaliação Justa:} A validação/teste utiliza obrigatoriamente $\beta=1.0$ fixo, garantindo comparação idônea entre os experimentos com e sem annealing (Teste A/B).
        \item \textbf{Resultado:} O modelo com annealing superou a baseline constante ao atingir menor ELBO de teste, mitigando o colapso do latente.
    \end{itemize}
\end{frame}

% ==============================================================================
% SLIDE 10: Resultados de Treinamento (Convergência)
% ==============================================================================
\begin{frame}{Resultados de Treinamento: Convergência}
    \begin{itemize}
        \item Comprovação empírica da queda estável da \textit{Loss} total (Negativa do ELBO).
        \item Equilíbrio entre o erro de reconstrução e a regularização após o \textit{Annealing}.
    \end{itemize}
    
    \begin{center}
        % ------------------------------------------------------------------
        % IMAGEM: vae_loss_curve.png
        % ------------------------------------------------------------------
        \includegraphics[width=0.75\textwidth]{vae_loss_curve.png}
        \\ \vspace{0.2cm}
        \footnotesize \textbf{Figura 2:} Curva de Perda do VAE monitorando treino e teste.
    \end{center}
    
    \vfill
\end{frame}

% ==============================================================================
% SLIDE 11: Resultados de Treinamento (Espaço Latente)
% ==============================================================================
\begin{frame}{Resultados de Treinamento: Organização Latente}
    \begin{itemize}
        \item Visualização das marginais 2D para as dimensões do espaço latente ($z \in \mathbb{R}^{32}$).
        \item Organização topológica densa e contínua, essencial para o estágio de difusão.
    \end{itemize}

    \begin{center}
        % ------------------------------------------------------------------
        % IMAGEM: latent_marginals.png
        % ------------------------------------------------------------------
        \includegraphics[width=0.85\textwidth]{latent_marginals.png}
        \\ \vspace{0.2cm}
        \footnotesize \textbf{Figura 3:} Marginais do espaço latente evidenciando aderência à normal $\mathcal{N}(0, I)$.
    \end{center}

    \vfill
\end{frame}

\section{3. Modelo de Difusão Latente}
% ==============================================================================
% SLIDE 12: Arquitetura LatentDenoiser
% ==============================================================================
\begin{frame}{Modelo de Difusão: O LatentDenoiser}
    Uma vez congelado o VAE ideal, o modelo de difusão é treinado exclusivamente operando no espaço latente de dimensionalidade reduzida ($m=32$).
    \vspace{0.2cm}
    \begin{itemize}
        \item \textbf{Objetivo do Denoiser:} Receber um vetor latente ruidoso $z_t$ e o indicador de tempo $t$, prevendo o ruído adicionado naquele passo para permitir a reversão determinística.
        \item \textbf{Hiperparâmetros de Treinamento:}
        \begin{itemize}
            \item \textbf{Batch Size:} 64
            \item \textbf{Learning Rate (LR):} 0.0001
            \item \textbf{Função de Perda:} Erro Quadrático Médio (\texttt{MSELoss}) sobre os ruídos previstos vs. reais.
            \item \textbf{Duração:} Treinado por 10 épocas, apresentando queda contínua e suave na curva de perda.
        \end{itemize}
    \end{itemize}
\end{frame}

% ==============================================================================
% SLIDE 13: Pipeline de Geração Completo
% ==============================================================================
\begin{frame}{Pipeline de Geração Completo}
    O fluxo final para conceber uma nova imagem sintética do EMNIST segue o algoritmo:
    \vspace{0.4cm}
    \begin{block}{Pipeline de Inferência Generativa}
        \centering
        Amostra de Ruído Puro $\mathcal{N}(0, I)$ \\
        $\Downarrow$ \\
        Processo Iterativo de \textbf{Denoising} (via \texttt{LatentDenoiser}) no Espaço Latente \\
        $\Downarrow$ \\
        Obtenção do Vetor Latente Limpo e Organizado $z_0 \in \mathbb{R}^{32}$ \\
        $\Downarrow$ \\
        Mapeamento Espacial pelo \textbf{VAE Decoder} \\
        $\Downarrow$ \\
        \textbf{Imagem Gerada Final} (Letra ou Dígito nítido no formato EMNIST)
    \end{block}
\end{frame}

\section{4. Conclusões}
% ==============================================================================
% SLIDE 14: Conclusões e Aprendizados
% ==============================================================================
\begin{frame}{Conclusões e Aprendizados}
    O desenvolvimento do projeto propiciou importantes constatações práticas no campo de IA Generativa:
    \vspace{0.2cm}
    \begin{itemize}
        \item \textbf{A importância do Acoplamento VAE-Difusão:} Modelos de difusão dependem crucialmente de que o espaço latente gerado pelo VAE seja contínuo e bem mapeado segundo a priori normal. O sucesso do \textbf{KL Annealing} foi o elo que viabilizou essa propriedade.
        \item \textbf{Sensibilidade à Likelihood:} Mudar de Gaussiana para Bernoulli alterou radicalmente a nitidez visual das amostras geradas, eliminando o efeito borrado crônico de VAEs básicas.
        \item \textbf{Eficiência Computacional:} Executar a difusão em $\mathbb{R}^{32}$ em vez de operar no espaço de pixels original economizou memória e reduziu o tempo de treino por época de forma drástica, validando a arquitetura de \textit{Latent Diffusion}.
    \end{itemize}
\end{frame}

\end{document}
